\documentclass[12pt]{article}

\usepackage[version=4]{mhchem}
\usepackage{chemfig}

\begin{document}
\section*{Vesiessee}
\bibliographystyle{ieeetr}
\renewcommand{\section}[2]{}

Vesi on yksi tärkeimmistä aineista maapallolla. Se koostuu kahdesta vetyatomista ja yhdestä happiatomista. Sen kemiallinen kaava on \ce{H2O}, ja
sen rakennekaava näyttää seuraavalta:
\\

\chemfig{O (-[:-127.75]H) (-[:-52.25]H)}
\\
\\

Vedessä kaikkien atomien välillä on kovalenttisia sidoksia, ja vesimolekyylien välillä on vetysidoksia.
Vedellä on poikkeuksellisen suuri kiehumis- ja sulamispiste. Tämä johtuu suoraan vesimolekyylien välillä olevista vetysidoksista.
Yllä olevaa rakennekaavaa tutkailemalla ilmenee että vesimolekyyli pystyy muodostamaan 4 eri vetysidosta muiden vesimolekyylin kanssa, sillä
molekyylillä on kaksi vetyatomia luovuttamassa ja happiatomin kautta kaksi vastaanottamassa.
Vetysidokset ovat suhteellisen voimakkaita heikkoja sidoksia, joten niiden rikkomiseen vaatii paljon lämpöenergiaa. On sanomatta selvää että jos niitä on neljä,
molekyylillä on suuri kiehumispiste.

Veden vetysidokset ovat merkittäviä myös muulla tavalla. Vesi on tiheimmillään noin 4 °C:n lämpötilassa. 
Kun vettä jäähdytetään huoneenlämpötilasta kohti 4 °C:n lämpötilaa, vesimolekyylien lämpöliike hidastuu ja molekyylit pääsevät pakkautumaan lähemmäs toisiaan. Tämän seurauksena veden tiheys kasvaa.

Alle 4 °C:n lämpötilassa vesimolekyylit alkavat kumminkin järjestyä muotoon, joka muistuttaa jään rakennetta. Tässä rakenteessa vesimolekyylien väliin jää enemmän tyhjää tilaa.

Tämän vuoksi vesi alkaa laajentua jäähtyessään alle 4 °C:n lämpötilaan ja sen tiheys pienenee. Jäässä tämä avoin rakenne on vielä voimakkaampi, minkä vuoksi jää on vettä harvempaa. Tästä myös johtuu miksi jää kelluu.
\\
\\
Vesi on poolinen aine. Tämän takia vesi liuottaa muita poolisia aineita. 
Tästä myös johtuu miksi ioniyhdisteet hajoaa veteen. Veden happiatomi on elektronegatiivisempi kuin vety, joten
happi on negatiivinen ja vety positiivinen. Kun vesimolekyyli tulee kontaktiin ioniyhdisteen kanssa, vesimolekyyli ja ioniyhditeet vetää toisiaan puoleen.
Vesi piirittäää ioniyhdisteen ionit erilleen toisistaan, ja pitää ne erillään.


\bibliography{testi.bib}

\end{document}